\subsection{RELATION WITH COXETER SYSTEMS}

THEOREM 1. Let \(\mathbf{C}\) be a chamber and let \(\mathbf{S}\) be the set of reflections with respect to the walls of \(\mathbf{C}\) .

\begin{enumerate}
\def\labelenumi{(\roman{enumi})}
\item
  The pair \((\mathbf{W},\mathbf{S})\) is a Coxeter system.
\item
  Let \(w \in \mathrm{W}\) and let \(\mathrm{H}\) be a wall of \(\mathrm{C}\) . The relation \(l(s_{\mathrm{H}}w) > l(w)\) implies that the chambers \(\mathrm{C}\) and \(w(\mathrm{C})\) are on the same side of \(\mathrm{H}\) .
\item
  For any chamber \(\mathbf{C}'\) , there exists a unique element \(w \in \mathbf{W}\) such that \(w(\mathbf{C}) = \mathbf{C}'\) .
\item
  The set of hyperplanes \(\mathbf{H}\) such that \(s_{\mathrm{H}} \in \mathbf{W}\) is equal to \(\mathfrak{H}\) .
\end{enumerate}

Every element of S is of order 2 and Lemma 2 shows that S generates W. For any wall H of C, denote by \(P_{H}\) the set of elements \(w \in W\) such that the chambers C and \(w(\mathrm{C})\) (which do not meet H) are on the same side of H. We shall verify conditions \((\mathrm{A}')\) , \((\mathrm{B}')\) and (C) of Chap. IV, §1, no. 7.

\((A') 1 \in P_H\) : Trivial.

(B') \(P_{H}\) and \(s_{H}.P_{H}\) are disjoint: Indeed, \(w(\mathrm{C})\) and \(s_{H}w(\mathrm{C})\) are on opposite sides of H, so if \(w(\mathrm{C})\) is on the same side of H as C, it is not on the same side as \(s_{\mathrm{H}}w(\mathrm{C})\) .

\begin{enumerate}
\def\labelenumi{(\Alph{enumi})}
\setcounter{enumi}{2}
\tightlist
\item
  Let \(w \in \mathrm{P_H}\) and let \(\mathrm{H}'\) be a wall of \(\mathbf{C}\) such that \(ws_{\mathrm{H}'} \notin \mathrm{P_H}\) ; then \(ws_{\mathrm{H}'} = s_{\mathrm{H}}w\) : By assumption, \(w(\mathbf{C})\) is on the same side of \(\mathbf{H}\) as \(\mathbf{C}\) and \(ws_{\mathrm{H}'}(\mathbf{C})\) is on the other side. Thus, \(ws_{\mathrm{H}'}(\mathbf{C})\) and \(w(\mathbf{C})\) are on opposite sides of H; hence, the chambers \(s_{\mathrm{H}^{\prime}}(\mathrm{C})\) and C are on opposite sides of the hyperplane \(w^{-1}(\mathrm{H})\) . Let a be a point of the face of C with support \(H^{\prime}\) . The point \(a = s_{\mathrm{H}^{\prime}}(a)\) is in the closure of the two chambers C and \(s_{\mathrm{H}^{\prime}}(\mathrm{C})\) that are contained respectively in the two open half-spaces bounded by \(w^{-1}(\mathrm{H})\) ; thus, \(a \in w^{-1}(\mathrm{H})\) , so \(\mathrm{H}^{\prime} = w^{-1}(\mathrm{H})\) . From this we deduce that \(s_{\mathrm{H}^{\prime}} = w^{-1}s_{\mathrm{H}}w\) , hence \(ws_{H^{\prime}} = s_{H}w\) .
\end{enumerate}

Assertions (i) and (ii) follow from this and Prop. 6 of Chap. IV, § 1, no. 7. Moreover, we have (loc. cit., condition (A))

\[
\bigcap_ {H \in \mathfrak {M}} P _ {H} = \{1 \}.\tag{3}
\]

Lemma 2 shows that W acts transitively on the set of chambers. Moreover, if \(w \in W\) is such that \(w(\mathrm{C}) = \mathrm{C}\) , then \(w \in P_{H}\) for every wall H of C, hence w = 1 by (3). This proves (iii).

Finally, let H be a hyperplane such that \(s_{H} \in W\) . If H did not belong to H, there would be at least one chamber \(C'\) meeting H ( \(\S1\) , no. 3, Prop. 7). Every point of \(H \cap C'\) is invariant under \(s_{H}\) , and thus belongs to the chambers \(C'\) and \(s_{\mathrm{H}}(\mathrm{C}')\) ; thus, \(\mathrm{C}' = s_{\mathrm{H}}(\mathrm{C}')\) , which contradicts (iii) since \(s_{H} \neq 1\) .

COROLLARY. Let \(\Sigma\) be a set of reflections generating W. Then every reflection belonging to W is conjugate to an element of \(\Sigma\) .

Let \(\mathfrak{H}'\) be the set of hyperplanes of the form \(w(\mathrm{H})\) with \(w \in \mathrm{W}\) and \(\mathrm{H} \in \mathfrak{H}\) such that \(s_{\mathrm{H}} \in \Sigma\) . Since \(\mathrm{W}\) is generated by the family \((s_{\mathrm{H}})_{\mathrm{H} \in \mathfrak{H}'}\) and since \(\mathfrak{H}'\) is stable under \(\mathrm{W}\) , we can apply all the results of this no. to \(\mathfrak{H}'\) instead of \(\mathfrak{H}\) ; but Th. 1, (iv) shows that every reflection in \(\mathrm{W}\) is of the form \(s_{\mathrm{H}}\) with \(\mathrm{H} \in \mathfrak{H}'\) , hence the corollary.

